\documentclass[../main.tex]{subfiles}
\begin{document}
%%%%%%%%%%%%%%%%%%%%%%%%%%%%%%%%%%%%%%%%%%%%%%%%%%%%%%%%%%%%%%%%%%%%%%%%%%%%%%%%%%%%%%%%%%%%%%%%%%%
\textbf{Copyright \copyright{} 2001-2021 The Omni Group. All rights reserved.}\\ OmniGraffle,
the OmniGraffle icon/logo, Graffle, Stenciltown, OmniOutliner, the OmniOutliner icon/logo, and
the Omni logo are registered trademarks of the Omni Group.

AirDrop, AirPrint, Apple, Apple Pencil, the Apple logo, Cocoa, Finder, iBooks, iPad, iPad Air,
iPad mini, iPad Pro, iPhone, iPod touch, macOS, Mac, the Mac logo, Multi-Touch, Safari, and
Touch Bar are trademarks of Apple Inc., registered in the U.S. and other countries. iBooks
Store is a registered service mark of Apple Inc., registered in the U.S. and other countries.

IOS is a trademark or registered trademark of Cisco in the U.S. and other countries and is used
under license.

Adobe and Photoshop are either registered trademarks or trademarks of Adobe Systems Incorporated
in the United States and/or other countries.

EPUB is a registered trademark of the IDPF.

\chapter*{Introduction}
\addcontentsline{toc}{chapter}{Introduction}

Meet OmniGraffle 7!

\ogfigure{0.800}{res/og7_intro_hero.png}

When development on OmniGraffle 7 started, we faced a veritable quandary in how we could make a
great---and beloved---app even better. Many of the new features you'll find in OmniGraffle are
inspired by feedback we have received from users like yourself. We listen. We take every bit
of input, whether from the outside or in, and use those requests as an opportunity to make
things better and easier for you to use.

Among the hundreds of updates we have made to OmniGraffle, our biggest new feature is
\textit{artboards}. Artboards are objects that reside on a special layer (an \textit{artboard
layer}) beneath your design, and are used to define export areas for objects on the layers
above:

\ogfigure{0.800}{res/og7_intro09.png}

To help you create artboards, we've added the new \textbf{Artboard Tool}, available in the Tool
Palette:

\ogfigure{0.800}{res/og7_intro10.png}

And to help facilitate exporting, we have completely redesigned the export experience, too:

\ogfigure{0.800}{res/og7_intro07.png}

After selecting a file type for the export, specify the export options, and then choose to
export your graphics at multiple resolutions: from 100\% scale (@1x), to 500\% scale (@5x), or
as large as 1000\% scale (@10x)!

\ogfigure{0.800}{res/og7_intro08.png}

Canvases, layers, and objects are easier to manage in a \textit{unified} Sidebar:

\ogfigure{0.800}{res/og7_intro01.png}

You can drill into layers and groups, change canvas, layer, and object names in the sidebar, and
change the size of the icons, too.

Use the Canvas Size inspector to choose a \textbf{Canvas Mode}:

\begin{itemize}
  \item \textbf{Fixed}: A canvas with its width and height set to a specific size; for example,
    1024 $\times$ 768 pixels.
  \item \textbf{Flexible}: A canvas which expands in a specified direction (up, down, left, or
    right) to suit your needs.
  \item \textbf{Infinite}: A canvas which extends infinitely in all directions.
\end{itemize}

\ogfigure{0.625}{res/og7_intro02.png}

The new \textbf{Point Editor} tool makes it easy for you to access a shape's vector points to
edit a shape and to work with Bézier curves.

\ogfigure{0.800}{res/og7_intro03.png}

\begin{ognote}
  Converting vector points to Bézier curves is a feature of OmniGraffle Pro.
\end{ognote}

And along those lines---pun intended---if you have OmniGraffle Pro, you can convert both text
and lines into shapes, and then apply any type of style or effect you can think of using the
Object inspectors.

\ogfigure{0.800}{res/og7_intro04.png}

\ogfigure{0.800}{res/og7_intro05.png}

If you use stencils, you can now choose where to place the Stencil Browser to better suit the
way you work. By default, the Stencil Browser is available in the right sidebar along with the
inspectors:

\ogfigure{0.800}{res/og7_stencils_rightsidebar.png}

However, you can configure OmniGraffle (via the General preferences pane) to make the Stencil
Browser available in the following locations:

\begin{itemize}
  \item \textbf{Left Sidebar}: For those times when you need to work with stencils and various
    inspectors at the same time, you can move the Stencil Browser to the left sidebar:

    \ogfigure{0.800}{res/og7_stencils_leftsidebar.png}
  \item \textbf{Popover}: Only occasionally use stencils? You can opt to display the Stencil
    Browser in a Popover Menu, available by clicking the \textbf{Stencils} button in the
    toolbar:

    \ogfigure{0.800}{res/og7_stencils_popover.png}
  \item \textbf{Floating Window}: If you have a larger display and tend to use stencils more
    frequently, you can choose to have the Stencil Browser available as a free-floating window
    that you can move around your display:

    \ogfigure{0.800}{res/og7_stencils_floatingwindow.png}
\end{itemize}

Need to use more than one stencil set at a time? Double-click a stencil set to view the stencils
in individual floating windows:

\ogfigure{0.800}{res/og7_stencils_multiplebrowsers.png}

Working with SVG? No problem! Just drag the SVG file into OmniGraffle to import the SVG source
and start working on your graphic. Or, choose SVG as the file type when exporting.

\ogfigure{0.800}{res/og7_intro11.png}

These features and many more are waiting for you in the new OmniGraffle 7---a thousand changes
in one radically new app for macOS.

\begin{ogtip}
  To see what's new in OmniGraffle 7, or to take a Walkthrough of OmniGraffle 7's key features,
  choose \textbf{Help} $\blacktriangleright$ \textbf{Welcome to OmniGraffle 7}.
\end{ogtip}

\section*{How This Guide is Organized}
\addcontentsline{toc}{section}{How This Guide is Organized}

The \textit{OmniGraffle 7 Reference Manual for macOS} that you have before you (either in
OmniGraffle itself, online, or in iBooks on your Mac or iOS device---or both) has been
assembled in such a way that you're provided with a complete tour of OmniGraffle.

This new edition is broken into three distinct parts, as follows:

\subsection*{Part I: OmniGraffle Basics}

The first part of the documentation introduces you to the basics of using OmniGraffle. From
finding your way around the interface, to customizing the user experience to suit your work,
these chapters hold the knowledge you need to start building OmniGraffle projects from
scratch:

\begin{itemize}
  \item Chapter 1, Basic Concepts --- This chapter explains some basic concepts of working in
    OmniGraffle, from working with the tools to draw things, inspectors to style things, and the
    sidebar to organize things.
  \item Chapter 2, Basic Setup and Configuration --- This chapter provides you with an in-depth
    look at the various ways you can customize OmniGraffle to provide you with a more
    \textit{personal} user experience.
  \item Chapter 3, Drawing Basics --- This chapter shows you how to create, connect, and style
    objects.
  \item Chapter 4, Diagramming Basics --- This short tutorial takes you through the process of
    creating a hierarchical diagram, while introducing you to the tools, inspectors, and other
    key elements of OmniGraffle.
\end{itemize}

\subsection*{Part II: OmniGraffle In-Depth}

Once you have a firm grounding in OmniGraffle, it's time to explore everything else that
OmniGraffle has to offer. This part introduces you to the new unified sidebar, artboards,
tools, inspectors, stencils, and exporting images from OmniGraffle.

\begin{itemize}
  \item Chapter 5, Managing Canvases, Layers, and Objects in the Sidebar --- Central to how you
    work in OmniGraffle, the new unified Sidebar makes it easy for you to manage canvases,
    layers, and all of the objects on those layers in one convenient place.
  \item Chapter 6, Using Artboards (PRO) --- Artboard Layers and Objects affect objects on the
    layers above, either acting as a container or for defining export areas in your projects.
  \item Chapter 7, Using the Tools --- The Toolbar contains all of the drawing tools you'll
    need, and this chapter provides you with a quick overview of what each of them do.
  \item Chapter 8, Working With the Inspectors --- Once you've learned about the Tools, you can
    use the Inspectors to examine the things you draw and modify their properties.
  \item Chapter 9, Using, Curating, and Creating Stencils --- This chapter shows you how to use
    and organize the stencils that come with OmniGraffle, and how to create new stencils for
    your own use or for sharing with other OmniGraffle users on
    \href{http://stenciltown.omnigroup.com}{Stenciltown}.
  \item Chapter 10, Using and Creating Custom Workspaces --- Workspaces enable you to define how
    you work with the inspectors in OmniGraffle. In addition to viewing the inspectors as a
    sidebar to the right of the canvas, or as a single or multiple windows for each inspector
    tab, you can create custom workspaces that define which inspectors are open or closed, and
    even have them appear in their own floating window.
  \item Chapter 11, Using OmniGraffle in Presentation Mode (PRO) --- Learn how to use
    OmniGraffle's Presentation Mode to present your work---designs and prototypes---on the big
    screen.
  \item Chapter 12, Using Templates and Diagram Styles --- This chapter provides you with an
    overview of Templates and Diagram Styles, and shows you how to create new ones for use in
    future projects.
  \item Chapter 13, Exporting Image and Graphics Files from OmniGraffle --- The new Export panel
    makes it easy for you to select the image or file output type, apply file type-specific
    options---and preview the images prior to exporting the files.
  \item Chapter 14, Using Data Variables --- In OmniGraffle, a variable is a snippet that you
    can enter as a label to an object or line, or as a text object on the canvas, which inserts
    other text. Learn how to use data variables to accelerate and make your work more dynamic.
\end{itemize}

\subsection*{Part III: Reference}

The back of the book consists of reference material; things that you may need on occasion, but
not frequently. Chapters include a complete listing of file types you can import or open in
OmniGraffle, the app's menu items (and their default keyboard shortcuts), how to sync files
with OmniPresence and Omni Sync Server, details for managing your OmniGraffle licenses, and
how to get help from our incredibly dedicated team of Support Humans.

\begin{itemize}
  \item Chapter 15, Opening and Importing Files in OmniGraffle --- One of the essential
    components of any great graphics application is its compatibility with various image and
    graphics filetypes.
  \item Chapter 16, Menus and Keyboard Shortcuts --- Here you'll find a detailed walkthrough of
    each menu item in OmniGraffle 7, including the Pro-only options. Where applicable, the
    descriptions include the menu option's keyboard shortcut so you can work even faster in
    OmniGraffle.
  \item Chapter 17, Getting Synced with OmniPresence and Omni Sync Server --- This chapter shows
    you how to download, install, and configure OmniPresence on your Mac so you can sync
    OmniGraffle files between your Mac and iOS devices.
  \item Chapter 18, Getting Help --- This chapter points you to all of the resources we have
    available: Support Humans, Support Articles, Videos, and Documentation.
\end{itemize}

\section*{What is Added in Pro?}
\addcontentsline{toc}{section}{What is Added in Pro?}

\begin{ogpro}
  This guide describes features common to both the Standard and Pro versions of OmniGraffle 7.
  Throughout the guide, look for boxes like this one for features that are Pro-only.

  Unsure which version you need? OmniGraffle 7 Pro includes advanced features for layers and
  drawing, working with text, exporting, displaying object details and notes, and for
  compatibility with Microsoft Visio\textsuperscript{\textregistered{}} and Adobe
  Photoshop\textsuperscript{\textregistered{}}.

  OmniGraffle 7 Pro features include:

  \begin{itemize}
    \item Omni Automation\ldots{}script OmniGraffle with JavaScript!
    \item Artboard Layers and Objects
    \item The Artboard tool
    \item Convert text and lines to shapes
    \item Text offsets and rotation
    \item Assign Actions to objects to create interactive prototypes
    \item AppleScript support
    \item Display object ID numbers in the Sidebar to aid in scripting
    \item Object fill blending
    \item Distortion effects
    \item Geometry left-center-right
    \item Vertex Point Geometry
    \item Define how objects appear on the canvas when at 100\% zoom
    \item Share layers between canvases
    \item Add and remove shape, line, and text styles to the Favorite Style section of the Style
      Well
    \item Combine shapes to create custom shapes and clipping paths
    \item Add Notes to and assign key-value pairs to objects
    \item Display Notes when mousing over objects on the canvas
    \item Create tables and subgraphs
    \item Work with multiple editing windows open
    \item Visio\textsuperscript{\textregistered{}} Support
    \item Presentation Mode
    \item Export layered OmniGraffle files as layered Adobe\textsuperscript{\textregistered{}}
      Photoshop\textsuperscript{\textregistered{}} files (.psd)
  \end{itemize}

  For a comparison between the Standard and Pro versions of OmniGraffle, see
  \href{https://www.omnigroup.com/omnigraffle\#pro}{the OmniGraffle product page} on our
  website.
\end{ogpro}
%%%%%%%%%%%%%%%%%%%%%%%%%%%%%%%%%%%%%%%%%%%%%%%%%%%%%%%%%%%%%%%%%%%%%%%%%%%%%%%%%%%%%%%%%%%%%%%%%%%
\end{document}
